% !TeX program = xelatex
\documentclass[a4paper,10pt]{article}

% ============================================================================
% ក្រដាសដាច់ដោយឡែក៖ សេណារីយ៉ូនៃការបង្រៀន (Scenario នៃការបង្រៀន – ៥០ នាទី)
% ============================================================================
\usepackage[a4paper, left=0.95cm, right=0.95cm, top=0.88cm, bottom=0.98cm, footskip=0.50cm]{geometry}
\usepackage{amsmath}
\usepackage{unicode-math}
\usepackage{fontspec}
\usepackage{array, longtable, colortbl}
\usepackage[table]{xcolor}
\usepackage{pifont}
\usepackage{fancyhdr}
\usepackage{tikz}

% LaTeX 2025 + longtable compatibility guard
\makeatletter
\let\@expl@@@mark@update@singlecol@structures@@\relax
\makeatother

% Cross-platform font path resolver (Windows local + Linux/Cloud repo fonts/)
\providecommand{\NGSFindFonts}{%
  \IfFileExists{C:/Windows/Fonts/cambria.ttc}{%
    \def\NGSSysPath{C:/Windows/Fonts/}%
    \def\NGSUserPath{C:/Users/KIT/AppData/Local/Microsoft/Windows/Fonts/}%
  }{%
    \IfFileExists{./fonts/KhmerOSsiemreap.ttf}{\def\NGSSysPath{./fonts/}\def\NGSUserPath{./fonts/}}{%
    \IfFileExists{../fonts/KhmerOSsiemreap.ttf}{\def\NGSSysPath{../fonts/}\def\NGSUserPath{../fonts/}}{%
    \IfFileExists{../../fonts/KhmerOSsiemreap.ttf}{\def\NGSSysPath{../../fonts/}\def\NGSUserPath{../../fonts/}}{%
    \IfFileExists{../../../fonts/KhmerOSsiemreap.ttf}{\def\NGSSysPath{../../../fonts/}\def\NGSUserPath{../../../fonts/}}{%
    \IfFileExists{../../../../fonts/KhmerOSsiemreap.ttf}{\def\NGSSysPath{../../../../fonts/}\def\NGSUserPath{../../../../fonts/}}{%
    \IfFileExists{../../../../../fonts/KhmerOSsiemreap.ttf}{\def\NGSSysPath{../../../../../fonts/}\def\NGSUserPath{../../../../../fonts/}}{%
    \IfFileExists{../../../../../../fonts/KhmerOSsiemreap.ttf}{\def\NGSSysPath{../../../../../../fonts/}\def\NGSUserPath{../../../../../../fonts/}}{%
      \def\NGSSysPath{../../../../../../../fonts/}\def\NGSUserPath{../../../../../../../fonts/}%
    }}}}}}}%
  }%
}
\NGSFindFonts

\setmathfont[
  Path={\NGSSysPath},
  FontIndex=1,
  Scale=1.02
]{cambria.ttc}

\defaultfontfeatures{Ligatures=TeX}

\setmainfont[
  Path={\NGSUserPath},
  Extension=.ttf,
  Script=Khmer,
  BoldFont={KhmerOSsiemreap},
  BoldFeatures={FakeBold=1.8},
  ItalicFont={KhmerOSsiemreap},
  ItalicFeatures={FakeSlant=0.25},
  BoldItalicFont={KhmerOSsiemreap},
  BoldItalicFeatures={FakeBold=1.8, FakeSlant=0.25}
]{KhmerOSsiemreap}

\newfontfamily\muol[
  Path={\NGSUserPath},
  Extension=.ttf,
  Script=Khmer
]{KhmerOSmuollight}

\newfontfamily\TimesEn[
  Path={\NGSSysPath},
  Extension=.ttf,
  Scale=1.00,
  BoldFont={timesbd},
  ItalicFont={timesi},
  BoldItalicFont={timesbi}
]{times}

\newfontfamily\NagaBold[
  Path={\NGSUserPath},
  Extension=.ttf
]{AKbalthom Naga Bold}

\newfontfamily\CalibriFont[
  Path={\NGSSysPath},
  Extension=.ttf
]{calibri}

\XeTeXlinebreaklocale "km"
\XeTeXlinebreakskip = 0pt plus 1pt minus 0.1pt

\definecolor{stepgray}{HTML}{D9D9D9}
\definecolor{headgray}{HTML}{F2F2F2}
\definecolor{darkred}{HTML}{C00000}
\definecolor{red5e}{HTML}{FF0000}

\pagestyle{fancy}
\fancyhf{}
\renewcommand{\headrulewidth}{0pt}
\fancyfoot[R]{\CalibriFont\fontsize{11pt}{13pt}\selectfont\thepage}

\setlength{\parindent}{0pt}
\setlength{\parskip}{0pt}
\setlength{\tabcolsep}{4pt}

\newcolumntype{L}[1]{>{\raggedright\arraybackslash}p{#1}}
\newcommand{\en}[1]{{\TimesEn #1}}
\newcommand{\enbf}[1]{{\TimesEn\textbf{#1}}}
\newcommand{\fiveE}[1]{{\NagaBold\fontsize{9pt}{12.5pt}\selectfont\textcolor{red5e}{#1}}}
\newcommand{\ditem}{\tikz[baseline=-0.5ex]\fill (0,0.40ex) -- (0.40ex,0) -- (0,-0.40ex) -- (-0.40ex,0) -- cycle;\hspace{0.25em}}
\newcommand{\chkbox}{\tikz[baseline=-0.1ex]\draw[black, line width=0.75pt](0,0) rectangle (0.23,0.23);}
\newcommand{\cnum}[1]{{\fontsize{9.5pt}{11pt}\selectfont\ding{\numexpr181+#1\relax}}\hspace{0.15em}}
\newcommand{\tri}[1]{\ensuremath{\triangle #1}}
\newcommand{\ang}[1]{\ensuremath{\angle #1}}
\newcommand{\cg}{\ensuremath{\cong}}
\newcommand{\ncg}{\ensuremath{\not\cong}}

\begin{document}
\fontsize{10pt}{14.8pt}\selectfont

% -------------------- ក្បាលឯកសារ Scenario --------------------
\begin{center}
  {\muol\fontsize{11pt}{16pt}\selectfont ព្រះរាជាណាចក្រកម្ពុជា}\\[1pt]
  {\muol\fontsize{10pt}{15pt}\selectfont ជាតិ សាសនា ព្រះមហាក្សត្រ}\\[4pt]
  {\muol\fontsize{10.5pt}{15pt}\selectfont សេណារីយ៉ូនៃការបង្រៀន \enbf{(Scenario} នៃការបង្រៀន\enbf{)}}
\end{center}

\vspace{2pt}
{\renewcommand{\arraystretch}{1.06}
\begin{tabular}{@{}l l l p{0.5cm} l l l@{}}
  - អនុវិទ្យាល័យ & : & ព្រះក្រសាំង (កំពង់ស្វាយ-កំពង់ធំ) & & - មុខវិជ្ជា & : & គណិតវិទ្យា ថ្នាក់ទី៨ \\
  - មេរៀនទី១២ & : & ១.៣.១ អ.ម (HA) និងលក្ខណៈពុះមុំ & & - រយៈពេលសរុប & : & \textcolor{darkred}{\textbf{១ ម៉ោងសិក្សា (៥០ នាទី)}} \\
  - គ្រូបង្រៀន & : & ហុក ប៊ុនថែន & & - កាលបរិច្ឆេទ & : & ថ្ងៃសៅរ៍ ទី២៨ ខែវិច្ឆិកា ឆ្នាំ២០២៦ \\
\end{tabular}
}

\vspace{3pt}

% -------------------- I. តារាងបែងចែកពេលវេលា និងសេណារីយ៉ូលម្អិត (៥០ នាទី) --------------------
\textbf{I.} {\muol តារាងបែងចែកពេលវេលា និងសេណារីយ៉ូប្រតិបត្តិការក្នុងថ្នាក់} \textbf{(៥០ នាទី)}

\renewcommand{\arraystretch}{1.10}
\begin{longtable}{|c|L{2.75cm}|c|L{7.65cm}|L{4.45cm}|}
\hline
\rowcolor{headgray}
\multicolumn{1}{|c|}{\muol នាទី} &
\multicolumn{1}{c|}{{\muol ជំហាន \enbf{/ 5Es}}} &
\multicolumn{1}{c|}{\muol រយៈពេល} &
\multicolumn{1}{c|}{{\muol សេណារីយ៉ូប្រតិបត្តិ \enbf{(}ពាក្យពេចន៍គ្រូ \enbf{\&} សម្ភារៈ\enbf{)}}} &
\multicolumn{1}{c|}{{\muol ការគ្រប់គ្រង \enbf{\&} វិធីបម្រុង}} \\
\hline
\endfirsthead

\hline
\rowcolor{headgray}
\multicolumn{1}{|c|}{\muol នាទី} &
\multicolumn{1}{c|}{{\muol ជំហាន \enbf{/ 5Es}}} &
\multicolumn{1}{c|}{\muol រយៈពេល} &
\multicolumn{1}{c|}{{\muol សេណារីយ៉ូប្រតិបត្តិ \enbf{(}ពាក្យពេចន៍គ្រូ \enbf{\&} សម្ភារៈ\enbf{)}}} &
\multicolumn{1}{c|}{{\muol ការគ្រប់គ្រង \enbf{\&} វិធីបម្រុង}} \\
\hline
\endhead

\hline
\endfoot

\hline
\endlastfoot

00--02 & \textbf{ជំហានទី១}\newline រដ្ឋបាលថ្នាក់ & ២ នាទី &
\ditem ពិនិត្យអវត្តមាន និងអនាម័យថ្នាក់រៀន ។ &
ចាត់ប្រធានក្រុមស្រាប់ពីមុន ។ \\[3pt]
& & &
\ditem ចាត់សិស្សជាក្រុម ៣ នាក់ («ក្រុមស្ថាបត្យករជណ្ដើរសុវត្ថិភាព») ដាក់កញ្ចប់សន្លឹក ១:១ (ទំព័រ ១) ទុកលើតុ ។ &
ហាមកាត់ក្រដាសមុនបញ្ជា ។ \\ \hline

02--05 & \textbf{ជំហានទី២ (ក)}\newline កំណែកិច្ចការផ្ទះ & ៣ នាទី &
\ditem \textbf{សំណួរបញ្ឆោត (\#4 \& \#9)}៖ «ហេតុអ្វី $\tri{ABC}\cg\tri{BDC}$ ខុស?» &
ផ្ដោតលើលំដាប់កំពូល $A\leftrightarrow D$ ។ \\[3pt]
& & &
\ditem ផ្ទៀងខ្លី \#4 ($\tri{ABD}\cg\tri{CBD}\to\tri{AED}\cg\tri{CED}$) និង \#9 ($\tri{ABC}\cg\tri{DCB}$ តាម ជ.ជ.ជ) ។ &
កុំហៅសិស្សសរសេរវែងលើក្ដារខៀន ។ \\ \hline

05--08 & \textbf{ជំហានទី២ (ខ)}\newline រំលឹកមេរៀនទាក់ទង & ៣ នាទី &
\ditem \textbf{ភាគ ១ (១.៥ ន.)}៖ បង្ហាញ $\tri{ABC}$ ($\ang{A}=90^\circ, \ang{C}=35^\circ$) ឱ្យស្គាល់អ៊ីប៉ូតេនុស $BC$ ឈម $90^\circ$ និង $\ang{B}=90^\circ{-}\ang{C}=55^\circ$ ។ &
បើសិស្សស្ទាក់ស្ទើររកអ៊ីប៉ូតេនុស ឱ្យគូសព្រួញចេញពីប្រអប់កែង $90^\circ$ ភ្លាម ។ \\[3pt]
& & &
\ditem \textbf{ភាគ ២ (១.៥ ន.)}៖ ប្រៀបធៀប $MA\perp OX$ ទល់នឹង $MD$ ទ្រេត ($\not\perp OX$) បញ្ជាក់ចម្ងាយវាស់តាម $MA\perp OX$ ។ &
រំលឹកពាក្យ «កែង = ខ្លីបំផុត» ។ \\ \hline

08--11 & \fiveE{E1 Engage}\newline ជណ្ដើរ \& ទ្វារទឹក & ៣ នាទី &
\ditem លើករឿង \textbf{«ជម្រាលជណ្ដើរសុវត្ថិភាព និងទ្វារទឹកប្រឡាយបំបែក»}៖ ជណ្ដើរ $5\text{m}$ មុំ $30^\circ$ ផ្អែកជញ្ជាំង $90^\circ$ + បង្គោល $M$ លើពុះមុំ $OM$ នៃប្រឡាយ $OX, OY$ តកែង $MA\perp OX, MB\perp OY$ ។ &
បង្ហាញរូប E1 លើក្ដារខៀន ឬស្លាយ ។ \\[3pt]
& & &
\ditem ស្រង់មតិទស្សន៍ទាយ រួចសរសេរចំណងជើង ១.៣.១ ។ &
កុំទាន់ផ្ដល់ចម្លើយផ្លូវការ ទុកជាបេសកកម្មឱ្យសិស្សស្រាយបញ្ជាក់ ។ \\ \hline

11--22 & \fiveE{E2 Explore}\newline ពិសោធន៍ ១ \& ២ & ១១ នាទី &
\ditem \textbf{សកម្មភាព ១ (៦ ន.)}៖ លើសន្លឹក ១:១ (ទំព័រ ១) ដាក់ខ្សែអ៊ីប៉ូតេនុស $5\text{cm}$ មុំ $35^\circ$ លើផ្ទាំងកែង $90^\circ$ កាត់ចេញ រួចបង្វិល-ត្រឡប់ត្រួតលើ $\tri{A'B'C'}$ ($\sim 135^\circ$) ស៊ីគ្នា ១០០\% ។ &
ត្រូវបោះពុម្ពទំព័រ ១ តាមខ្នាត \textbf{Actual Size 100\%} ។ \\[3pt]
& & &
\ditem \textbf{សកម្មភាព ២ POE 2A (២.៥ ន.)}៖ ទស្សន៍ទាយ រួចត្រួត $\tri{PQR}$ (អ៊ីប៉ូតេនុស $5\text{cm}, 35^\circ$) លើ $\tri{P'Q'R'}$ (ជ្រុងកែង $5\text{cm}, 35^\circ$, អ៊ីប៉ូតេនុស $\approx 6.1\text{cm}$) $\Rightarrow \ncg$ &
ឱ្យសិស្សទស្សន៍ទាយមុនត្រួតក្រដាស ។ \\[3pt]
& & &
\ditem \textbf{សកម្មភាព ២ POE 2B (២.៥ ន.)}៖ បត់ក្រដាស $\ang{XOY}$ រកផ្នត់ពុះមុំ $OM$ និងកែង $MA\perp OX, MB\perp OY$ ($MA=MB$) ធៀបនឹងទ្រេត $MD\not\perp OY$ ($MD > MB$) ។ &
គ្រូដើរផ្ទៀងការបត់ $MA\perp OX$ ឱ្យត្រួតស្មើជ្រុង $OX$ ។ \\ \hline

22--28 & \fiveE{E3 Explain}\newline សម្រាយ \& \cnum{1}--\cnum{5} & ៦ នាទី &
\ditem ដឹកនាំសម្រាយ $90^\circ{-}\ang{C}\to$ ម.ជ.ម និង\textbf{កែសៀវភៅទំព័រ ១៥៣ ($\ang{A}=\ang{A'}\to\ang{B}=\ang{B'}$)} ។ &
ចង្អុលកំហុស $\ang{A}=\ang{A'}$ ក្នុងសៀវភៅទំព័រ ១៥៣ ឱ្យច្បាស់ ។ \\[3pt]
& & &
\ditem បង្គុំទ្រឹស្ដីបទ \textbf{អ.ម (HA)} លើរូបបង្វិល $\sim 135^\circ$ + បង្រៀន \textbf{ប្រអប់កែង $90^\circ\to$ ព្រួញចង្អុលអ៊ីប៉ូតេនុស} ។ &
គូសរូបទ្រឹស្ដីបទបង្វិល $\sim 135^\circ$ កុំគូសតែរូបឈរត្រង់ស្ដង់ដារ ។ \\[3pt]
& & &
\ditem បិទផ្ទាំង \textbf{ធនាគារហេតុផល \cnum{1}--\cnum{5}} (ទំព័រ ៤) ។ & \\ \hline

28--38 & \fiveE{E4 Elaborate}\newline ព្រួញ $OM$ \& ពុះមុំ & ១០ នាទី &
\ditem ចែកកាត \textbf{«កាតព្រួញឈម $90^\circ$ និងអ៊ីប៉ូតេនុសរួម $OM$»} (ទំព័រ ២, ទំ.១៥៣)៖\newline
+ ២.៥ន. \textbf{កែរូបទំ.១៥៣ (ដក $/$ ពី $OA, OB$)} + កែពាក្យ \textbf{$OMA$ និង $OMB$} រួច \textbf{[tracing]} $90^\circ\to OM$ &
សង្កត់ធ្ងន់មុនសិស្សចាប់ខ្មៅដៃ៖ ហាមសរសេរ $OA=OB$ ជាបម្រាប់! \\[3pt]
& & &
+ ៥ន. សរសេរសម្រាយ \cnum{1}--\cnum{5} ($MA=MB$) និងបញ្ចប់ប្រអប់ទ្រឹស្ដីបទពុះមុំផ្ទាល់-ច្រាស (ទំ.១៥៤, ប្រកាសច្រាសស្រាយនៅ អ.ជ / HL) &
ព្រួញពីកែង $A, B$ ចង្អុលទៅ $OM$ ដូចគ្នា ។ \\[3pt]
& & &
+ ២.៥ន. \textbf{EIA 30s} រំកិល $M$ តាម $OM$ ញែកអ្វីប្រែ vs ថេរ ។ &
សង្ខេប៖ $OM, MA, MB$ ប្រែ តែ $MA=MB$ ថេរ ។ \\ \hline

38--48 & \fiveE{E5 Evaluate}\newline ស្ទូឌីយោ \& កាត A--D & ១០ នាទី &
\ditem \textbf{«ស្ទូឌីយោផ្ទៀងកម្ពស់ខ្វែង ៣ តួនាទី» (៧ ន.)} (ទំព័រ ៣, ប្រតិបត្តិទំ.១៥៤ $\to$ បំបែក $\tri{ADB}, \tri{AEC}$)៖\newline
+ \textbf{តួនាទី ១ (Pause-Ask)}៖ ពន្យល់មុនគូស $90^\circ\to$ អ៊ីប៉ូតេនុស $AB, AC$ &
គ្រូដើរត្រួតរូបបំបែក $\tri{ADB}, \tri{AEC}$ តាមក្រុម ។ \\[3pt]
& & &
+ \textbf{តួនាទី ២}៖ សរសេរ \cnum{1}--\cnum{4} ($\tri{ADB}\cg\tri{AEC}$ តាម អ.ម)\newline
+ \textbf{តួនាទី ៣ (Pause-Review)}៖ ផ្ទៀង $BD=CE$ (\cnum{5}) ដក $BE=AB{-}AE=AC{-}AD=CD$ + ទទួលស្គាល់ចម្លើយទី២ ($\tri{BDC}\cg\tri{CEB}$) + \textbf{បិទ E1} ។ &
ទទួលស្គាល់ទាំងវិធី $\tri{ADB}\cg\tri{AEC}$ និង $\tri{BDC}\cg\tri{CEB}$ ។ \\[3pt]
& & &
\ditem \textbf{កាតត្រួតពិនិត្យ A--D (៣ ន. / 60--90s)} (ទំព័រ ៥)៖ (A) \ding{51} អ.ម $135^\circ$, (B) \ding{55} អ៊ីប៉ូ vs ជ្រុង, (C) \ding{55} ទ្រេត $\not\perp$, (D) \ding{55} AA គ្មានជ្រុង + ម្រាម ១--៣ ។ &
បើតួនាទី ៣ អស់ពេលយូរ កាត់កាត A--D មក ៩០ វិនាទីរហ័ស ដើម្បីចប់ទាន់ ៥០ នាទី ។ \\ \hline

48--50 & \textbf{ជំហានទី៥}\newline កិច្ចការផ្ទះ \& បណ្ដាំ & ២ នាទី &
\ditem ចែកប័ណ្ណកិច្ចការផ្ទះ (ទំព័រ ៦)៖ ទី៦ (ទំ.១៦៤ ប្រយ័ត្ន $AD$ ជាជ្រុងកែង $\to$ ម.ជ.ម) \& ទី៧ (ទំ.១៦៥ ក. $PB=PC$ តាម ជ.ម.ជ, ខ. $PH\perp AB, PK\perp AC \Rightarrow PH=PK$ ដោយ អ.ម) ។ &
បញ្ជូនរូបលំហាត់ចូល Telegram / Classroom ជាស្រេច ។ \\[3pt]
& & &
\ditem ផ្ដាំផ្ញើសុវត្ថិភាពចរាចរណ៍ និងជួយការងារផ្ទះ ។ &
បញ្ចប់ម៉ោងសិក្សាឱ្យទាន់ ៥០ នាទី ។ \\ \hline

\rowcolor{stepgray}
\multicolumn{2}{|r|}{\textbf{សរុបរយៈពេលបង្រៀនពេញ}} & \textbf{៥០ នាទី} &
\multicolumn{2}{l|}{ត្រូវគ្នា ១០០\% នឹង \en{LS\_Plan\_HA.tex} និង \en{Printable\_Materials\_HA.tex} (ទំព័រ ១--៦)} \\ \hline
\end{longtable}

\vspace{1pt}

% -------------------- II. បញ្ជីត្រៀមសម្ភារៈ & III. ការគ្រប់គ្រងកំហុស --------------------
{\fontsize{9.2pt}{13.0pt}\selectfont\renewcommand{\arraystretch}{1.04}
\begin{minipage}[t]{0.44\textwidth}
  \textbf{II.} {\muol បញ្ជីត្រៀមសម្ភារៈមុនចូលថ្នាក់} \textbf{(ស-ខ្មៅ)}\\[2pt]
  \begin{tabular}{|c|L{6.75cm}|}
    \hline
    \rowcolor{headgray}
    \textbf{\ding{51}} & \multicolumn{1}{c|}{\textbf{សម្ភារៈត្រូវរៀបចំទុកជាមុន}} \\ \hline
    \chkbox & សន្លឹក ១:១ ទំព័រ ១ (\textbf{Actual Size 100\%})៖ ផ្ទាំងកែង $90^\circ$, បន្ទះ $5\text{cm}, 35^\circ$, រូប POE 2A \& ក្រដាសបត់ POE 2B \\ \hline
    \chkbox & កន្ត្រៃ កែង $90^\circ$ រ៉ាព័រទ័រ និងបន្ទាត់ក្រិតសង់/ផ្ទៀង ១:១ \\ \hline
    \chkbox & កាត E4 (ទំព័រ ២), សន្លឹក ៣ តួនាទី E5 (ទំព័រ ៣) \& ផ្ទាំងធនាគារ \cnum{1}--\cnum{5} (ទំព័រ ៤) \\ \hline
    \chkbox & កាត Exit \textbf{A--D} (ទំព័រ ៥) \& ប័ណ្ណកិច្ចការផ្ទះទី៦, ទី៧ (ទំព័រ ៦) \\ \hline
  \end{tabular}
\end{minipage}%
\hfill
\begin{minipage}[t]{0.54\textwidth}
  \textbf{III.} {\muol កំហុសភ័ន្តច្រឡំ} \textbf{(តាមការស្រាវជ្រាវ) \&} {\muol ដំណោះស្រាយ}\\[2pt]
  \begin{tabular}{|L{9.25cm}|}
    \hline
    \rowcolor{headgray}
    \multicolumn{1}{|c|}{\textbf{ឧបសគ្គសិស្ស (Empirical Misconceptions) \& វិធានការទប់ស្កាត់}} \\ \hline
    \textbf{១. កម្ពស់កែង (8\% មុនរៀន) \& ទិសដៅបង្វិល (ធ្លាក់ដល់ 21.5\%)៖} ប្រើ \textit{ប្រអប់កែង $90^\circ \to$ ព្រួញចង្អុលអ៊ីប៉ូតេនុស} និងបង្វិលកាត់ $\sim 135^\circ$ \en{(Yıldız et al., 2010; Ubah \& Bansilal, 2019)} ។ \\ \hline
    \textbf{២. អ៊ីប៉ូតេនុស vs ជ្រុងកែង \& ចម្ងាយទ្រេត $\not\perp$៖} ប្រើ POE 2A ($5\text{cm}$ vs $\approx 6.1\text{cm}$) និងបត់ក្រដាស POE 2B ($MA\perp OX$ vs $MD$ ទ្រេត) \en{(Šipuš et al., 2022; Catman-Aksoy \& Isiksal-Bostan, 2024)} ។ \\ \hline
    \textbf{៣. ត្រីកោណត្រួតគ្នា (92\% $\to$ 53\%) \& គ្មានរូប (72.5\% gaps)៖} គូសបំបែក $\tri{ADB}, \tri{AEC}$ (ទំព័រ ៣) ប្រើ \textbf{[tracing]} \& ធនាគារ \cnum{1}--\cnum{5} \en{(Biber, 2020; Winer \& Battista, 2022)} ។ \\ \hline
    \textbf{៤. កំហុសសៀវភៅទំ.១៥៣ \& អន្ទាក់ \textbf{/} លើ $OA, OB$៖} កែ $\ang{A}=\ang{A'}\to\ang{B}=\ang{B'}$, កែ $OMA$ និង $OMB$, លុប $/$ ចេញពី $OA, OB$ និងបំបែកផ្ទាល់ (HA) ពីច្រាស (HL) \en{(Shahbari \& Daher, 2020)} ។ \\ \hline
  \end{tabular}
\end{minipage}
}

\end{document}
